\section*{Capítulo \ref{ch:g12:reaction-rates} --- A velocidade das reações}

\begin{solution}{exo:g12:reaction-rates:1}
Temperatura; área de contato; concentração; temperatura.
\end{solution}

\begin{solution}{exo:g12:reaction-rates:2}
\qty{13.9}{min} (azul) e \qty{6.9}{min} (vermelha). Depois de duas
meias-vidas, $10.0 / 4 = \qty{2.5}{mmol/L}$.
\end{solution}

\begin{solution}{exo:g12:reaction-rates:3}
Inclinação $(5.0 - 1.0)/10 = \qty{0.40}{mmol.L^{-1}.min^{-1}}$: essa é a
\omterm{def:g12:reaction-rates:rate}{velocidade de formação} em \qty{5}{min}.
\end{solution}

\begin{solution}{exo:g12:reaction-rates:4}
O resfriamento e a \omterm{def:g10:concentration-and-dilution:dilution}{diluição} tornam a reação tão lenta que a composição
não muda mais durante a \omterm{def:g11:titration:titration}{titulação}. A composição medida é a que a \omterm{def:g6:pure-substances-and-mixtures:mixture}{mistura}
tinha quando a reação foi interrompida, no momento em que foi despejada.
\end{solution}

\begin{solution}{exo:g12:reaction-rates:5}
\qty{30}{min} são 3 meias-vidas: resta $1/8$. \qty{1}{h} são 6
meias-vidas: $1/64$, cerca de \qty{1.6}{\percent}.
\end{solution}

\begin{solution}{exo:g12:reaction-rates:6}
$\ln[\ce{A}]$: 2.079, 1.733, 1.386, 1.040, 0.693. Ele diminui 0.346 a
cada \qty{5}{min}: uma reta, primeira ordem. $k = 0.346/5 =
\qty{0.0693}{min^{-1}}$ e $t_{1/2} = \ln 2 / k = \qty{10.0}{min}$ (como
se vê diretamente: 8.00, 4.00, 2.00 a cada \qty{10}{min}).
\end{solution}

\begin{solution}{exo:g12:reaction-rates:7}
$k = 0.693 / 25 = \qty{0.0277}{s^{-1}}$.
\end{solution}

\begin{solution}{exo:g12:reaction-rates:8}
$0.100\,e^{-0.20} = \qty{0.0819}{mol/L}$; $0.100\,e^{-0.60} =
\qty{0.0549}{mol/L}$; $0.100\,e^{-2.0} = \qty{0.0135}{mol/L}$.
\end{solution}

\begin{solution}{exo:g12:reaction-rates:9}
Em \qty{13.9}{min}, $[\ce{A}] = \qty{5.0}{mmol/L}$; a inclinação da
tangente é próxima de $-0.25$, e $-k[\ce{A}] = -0.050 \times 5.0 =
\qty{-0.25}{mmol.L^{-1}.min^{-1}}$: a metade do valor inicial.
\end{solution}

\begin{solution}{exo:g12:reaction-rates:10}
O iodo é a única espécie colorida: pela lei de Beer--Lambert,
$A = \varepsilon \ell [\ce{I2}]$. A \omterm{def:g12:reaction-rates:rate}{velocidade de formação} é a
inclinação da tangente à curva $[\ce{I2}](t)$, isto é, a inclinação da
tangente a $A(t)$ dividida por $\varepsilon \ell$.
\end{solution}

\begin{solution}{exo:g12:reaction-rates:11}
A velocidade inicial $k[\ce{A}]_0$ é duas vezes maior no segundo
experimento; as meias-vidas são iguais, $\ln 2 / k$, independentes de
$[\ce{A}]_0$.
\end{solution}

\begin{solution}{exo:g12:reaction-rates:12}
$t = \ln(1/f)/k$. Em meias-vidas: $\ln(1/f)/\ln 2$. Para
\qty{1}{\percent}, $\ln 100 / \ln 2 = 6.6$; para \qty{0.1}{\percent},
$\ln 1000 / \ln 2 =
10.0$.
\end{solution}

\begin{solution}{exo:g12:reaction-rates:13}
\qty{90}{\percent} consumidos, $f = 0.10$: $t = \ln 10 / k$, isto é,
\qty{46}{min} a \qty{20}{\celsius} e \qty{23}{min} a \qty{30}{\celsius}.
A \omterm{def:g12:reaction-rates:first-order}{constante de velocidade} dobra em \qty{10}{\celsius}: a regra vale
aqui.
\end{solution}

\begin{solution}{exo:g12:reaction-rates:14}
A temperatura cai para perto de \qty{0}{\celsius}, e as concentrações
são divididas por 10. Com uma velocidade proporcional ao produto de duas
concentrações, a \omterm{def:g10:concentration-and-dilution:dilution}{diluição} sozinha a divide por $10 \times 10 = 100$.
\end{solution}

\begin{solution}{exo:g12:reaction-rates:15}
$v(t) = -\dfrac{d}{dt}\big([\ce{A}]_0 e^{-kt}\big) = k[\ce{A}]_0 e^{-kt}$.
Velocidade inicial $k[\ce{A}]_0$; a velocidade cai à metade quando
$e^{-kt} = 1/2$, isto é, depois de uma \omterm{def:g12:reaction-rates:half-life}{meia-vida}.
\end{solution}

\begin{solution}{pb:g12:reaction-rates:1}
\textbf{1.} Pela lei de Beer--Lambert, sendo o corante a única espécie
que absorve nesse comprimento de onda; a solução deve ser diluída o
bastante ($A$ abaixo de cerca de 1).

\textbf{2.} Ele não absorve nenhuma luz visível: não acrescenta nada a
$A$.

\textbf{3.} Para que a sua concentração quase não mude: a velocidade
passa então a depender só do corante.

\textbf{4.} $A = 0$ (o branco).

\textbf{5.} $A$ cai de 0.800 para 0.402 em \qty{6}{min}: cerca de
\qty{6}{min}.

\textbf{6.} De 0.402 em \qty{6}{min} para 0.199 em \qty{12}{min}: de
novo dividida por dois em \qty{6}{min}.

\textbf{7.} $-0.223$, $-0.448$, $-0.691$, $-0.911$, $-1.155$, $-1.366$,
$-1.614$, $-1.945$, $-2.551$, $-3.079$.

\textbf{8.} Os pontos ficam numa reta: a reação é de primeira ordem em
relação ao corante.

\textbf{9.} $(-2.551 + 0.223)/20 = \qty{-0.116}{min^{-1}}$.

\textbf{10.} $k = \qty{0.116}{min^{-1}}$.

\textbf{11.} $0.693 / 0.116 = \qty{6.0}{min}$, como lido na tabela.

\textbf{12.} $k A_0 = 0.116 \times 0.800 = \qty{0.093}{min^{-1}}$, em
unidades de \omterm{def:g11:absorbance:absorbance}{absorbância} por minuto.

\textbf{13.} A mesma, \qty{6.0}{min}: a \omterm{def:g12:reaction-rates:half-life}{meia-vida} de uma \omterm{def:g12:reaction-rates:first-order}{reação de
primeira ordem} não depende do valor inicial.

\textbf{14.} $0.800\, e^{-0.116 \times 30} = 0.025$.

\textbf{15.} $t = \ln 100 / k = 4.61 / 0.116 = \qty{40}{min}$.

\textbf{16.} $40 / 6.0 = 6.6$ meias-vidas.

\textbf{17.} $t = \ln(0.800/0.010)/k = 4.38 / 0.116 = \qty{38}{min}$.

\textbf{18.} Ainda 6.6 meias-vidas, agora
$6.6 \times 3.0 = \qty{20}{min}$.

\textbf{19.} Cerca de 6.6 meias-vidas, isto é, cerca de \qty{40}{min}
aqui.
\end{solution}
